% ============================================================
%  CHAPTER 14: GAS LAWS AND KINETIC THEORY
%  The Physics Codex: A Complete University Foundation
%  Korede Samuel Omotosho (Falcon)
%  Aurikrex Learning Systems, 2026
%
%  File: chapters/part2/ch14_gas_laws.tex
% ============================================================

\chapter{Gas Laws and Kinetic Theory}
\label{ch:gas_laws}
\minitoc
\newpage

\chapterquote{The atoms come into my brain, dance a
dance, and then go out --- there are always new atoms,
but always doing the same dance.}{Richard Feynman}

\begin{objectivebox}
By the end of this chapter, you should be able to:
\begin{enumerate}[noitemsep]
  \item State and apply Boyle's Law, Charles' Law
  and the Pressure Law
  \item Combine the three laws into the general gas
  equation
  \item State and apply the ideal gas equation
  $PV = nRT$
  \item Explain the assumptions of the kinetic theory
  of gases
  \item Describe qualitatively how kinetic theory
  explains the gas laws
  \item Connect average molecular kinetic energy to
  temperature: $\frac{1}{2}m\langle c^2 \rangle
  = \frac{3}{2}k_BT$
  \item Distinguish between ideal and real gases
  \item Solve problems involving mixtures of gases
  using Dalton's Law
\end{enumerate}
\end{objectivebox}

% ============================================================
\section{Intuition: What is a Gas?}
% ============================================================

Pump air into a bicycle tyre and it becomes harder ---
more pressure. Heat a sealed container of gas and the
pressure rises. Compress a gas and its temperature
rises. These are everyday observations, but they point
to deep physics.

A gas consists of an enormous number of molecules
moving randomly in all directions at high speeds.
At room temperature, nitrogen molecules (which make
up 78\% of air) move at average speeds of around
$500\ \text{m\,s}^{-1}$ --- faster than a bullet.
They collide with each other and with the walls of
their container billions of times per second. Each
collision with a wall exerts a tiny force. The sum
of billions of these tiny forces over the wall area
is what we measure as \textbf{pressure}.

This picture --- the \textbf{kinetic theory of gases}
--- connects the microscopic (invisible molecular
motion) to the macroscopic (measurable pressure,
volume, temperature). It is one of the great
intellectual achievements of 19th-century physics.

% ============================================================
\section{Boyle's Law}
% ============================================================

\begin{historybox}[Robert Boyle, 1662]
Robert Boyle (1627--1691), working in Oxford with
his assistant Robert Hooke, used a J-shaped glass
tube sealed at one end and open at the other, with
mercury trapping a column of air. By adding mercury
to the open end, he increased the pressure on the
trapped air and measured how its volume changed.
His meticulous measurements, published in 1662,
established the first quantitative gas law. Boyle
was a founding member of the Royal Society and is
often called the ``father of modern chemistry.''
\end{historybox}

\begin{lawbox}[Boyle's Law]
At constant temperature, the volume of a fixed mass
of gas is inversely proportional to its pressure:
\[
  P \propto \frac{1}{V}
  \quad \text{(constant } T \text{ and mass)}
\]
\[
  PV = \text{constant}
\]
\[
  \boxed{P_1 V_1 = P_2 V_2}
  \quad \text{(at constant } T\text{)}
\]
\end{lawbox}

\begin{diagbox}[Boyle's Law: $P$ vs $V$ and $P$ vs $1/V$ Graphs]
\begin{center}
\begin{tikzpicture}
  % P vs V graph
  \begin{axis}[
    name=plot1,
    width=6cm, height=5.5cm,
    xlabel={Volume $V$},
    ylabel={Pressure $P$},
    xmin=0, xmax=5,
    ymin=0, ymax=5,
    grid=major, grid style={gray!15},
    xtick=\empty, ytick=\empty,
    title={$P$ vs $V$ (hyperbola)}
  ]
    \addplot[codexblue, very thick, domain=0.3:4.8,
      samples=80]{1.0/x + 0.3};
    \addplot[codexred, very thick, domain=0.5:4.8,
      samples=80]{2.0/x + 0.3};
    \addplot[codexgreen, very thick, domain=0.8:4.8,
      samples=80]{3.5/x + 0.3};
    \node[codexblue, font=\tiny] at (axis cs:0.7,2.5)
      {$T_1$};
    \node[codexred, font=\tiny] at (axis cs:0.9,3.5)
      {$T_2 > T_1$};
    \node[codexgreen, font=\tiny] at (axis cs:1.5,4.0)
      {$T_3 > T_2$};
    \node[font=\tiny, gray] at (axis cs:3.5,1.5)
      {\shortstack{$PV = $ const\\(each curve)}};
  \end{axis}

  % P vs 1/V graph
  \begin{axis}[
    name=plot2,
    at={(plot1.right of south east)},
    anchor=left of south west,
    xshift=0.8cm,
    width=6cm, height=5.5cm,
    xlabel={$1/V$},
    ylabel={Pressure $P$},
    xmin=0, xmax=5,
    ymin=0, ymax=5,
    grid=major, grid style={gray!15},
    xtick=\empty, ytick=\empty,
    title={$P$ vs $1/V$ (straight line through origin)}
  ]
    \addplot[codexblue, very thick, domain=0:4.5]{0.8*x};
    \addplot[codexred, very thick, domain=0:3.5]{1.2*x};
    \addplot[codexgreen, very thick, domain=0:2.8]{1.6*x};
    \node[codexblue, font=\tiny] at (axis cs:4.0,2.8)
      {$T_1$};
    \node[codexred, font=\tiny] at (axis cs:3.0,3.5)
      {$T_2$};
    \node[codexgreen, font=\tiny] at (axis cs:2.0,3.5)
      {$T_3$};
    \node[font=\tiny, gray, align=center]
      at (axis cs:3.0,1.0)
      {Gradient $= $ const\\$= $ slope $\propto T$};
  \end{axis}
\end{tikzpicture}
\end{center}
\small Each curve on the $P$-$V$ graph corresponds to
a different temperature (isotherm). The $P$ vs $1/V$
graph is a straight line through the origin,
confirming $P \propto 1/V$.
\end{diagbox}

\begin{examplebox}[14.1]
\stepcounter{workedexample}
A gas occupies $4.0\ \text{L}$ at a pressure of
$200\ \text{kPa}$ at constant temperature.\\
(a) Find the pressure when the volume is compressed
to $1.5\ \text{L}$.\\
(b) Find the volume when the pressure is
$50\ \text{kPa}$.
\end{examplebox}

\begin{solutionbox}
\textbf{(a)} At constant $T$: $P_1V_1 = P_2V_2$
\[
  P_2 = \frac{P_1V_1}{V_2}
  = \frac{200 \times 4.0}{1.5}
  = 533\ \text{kPa}
\]

\textbf{(b)}
\[
  V_2 = \frac{P_1V_1}{P_2}
  = \frac{200 \times 4.0}{50}
  = 16.0\ \text{L}
\]
\end{solutionbox}

% ============================================================
\section{Charles' Law}
% ============================================================

\begin{historybox}[Jacques Charles and Gay-Lussac, 1787--1802]
Jacques Charles (1746--1823), a French physicist and
balloonist, discovered in 1787 that different gases
expand equally when heated through the same temperature
interval --- but he did not publish his findings.
Joseph Gay-Lussac rediscovered and published the law
in 1802, crediting Charles for the earlier work.
Charles was a passionate hot-air balloonist --- the
first person to use hydrogen balloons for flight ---
so he had a very practical interest in how gases
behave when heated.
\end{historybox}

\begin{lawbox}[Charles' Law]
At constant pressure, the volume of a fixed mass of
gas is directly proportional to its absolute temperature
(in Kelvin):
\[
  V \propto T
  \quad \text{(constant } P \text{ and mass)}
\]
\[
  \frac{V}{T} = \text{constant}
\]
\[
  \boxed{\frac{V_1}{T_1} = \frac{V_2}{T_2}}
  \quad \text{(at constant } P\text{)}
\]
\textbf{Critical:} $T$ must be in \textbf{Kelvin}.
Using Celsius gives wrong answers.
\end{lawbox}

\begin{diagbox}[Charles' Law: $V$ vs $T$ Graph]
\begin{center}
\begin{tikzpicture}
\begin{axis}[
  width=10.5cm, height=6cm,
  xlabel={Temperature $T$ (K)},
  ylabel={Volume $V$},
  xmin=-300, xmax=600,
  ymin=0, ymax=5,
  grid=major, grid style={gray!15},
  xtick={-273,0,200,400},
  xticklabels={\shortstack{$0\,\text{K}$\\$-273\textdegree\text{C}$},$0\textdegree\text{C}$,$200\textdegree\text{C}$,$400\textdegree\text{C}$},
  ytick=\empty,
  extra x ticks={-273},
  extra x tick labels={},
]
  % Three lines for three different pressures
  \addplot[codexblue, very thick, domain=-273:550]
    {0.008*(x+273)};
  \addplot[codexred, very thick, domain=-273:550]
    {0.012*(x+273)};
  \addplot[codexgreen, very thick, domain=-273:550]
    {0.006*(x+273)};

  % Extrapolation to absolute zero (dashed)
  \addplot[codexblue, dashed, domain=-273:0]
    {0.008*(x+273)};
  \addplot[codexred, dashed, domain=-273:0]
    {0.012*(x+273)};
  \addplot[codexgreen, dashed, domain=-273:0]
    {0.006*(x+273)};

  % All lines converge at absolute zero
  \fill[codexred] (axis cs:-273,0) circle (4pt);
  \node[below right, codexred, font=\small]
    at (axis cs:-250,0.35) {\shortstack{Absolute zero\\$T = 0\,\text{K}$}};

  % Pressure labels
  \node[codexblue, font=\small]
    at (axis cs:480,3.2) {$P_1$};
  \node[codexred, font=\small]
    at (axis cs:400,4.5) {$P_2 > P_1$};
  \node[codexgreen, font=\small]
    at (axis cs:480,2.3) {$P_3 < P_1$};

  % Real gas region
  \draw[gray, dashed, thick]
    (axis cs:0,-0.1) -- (axis cs:0,5);
  \node[gray, font=\tiny, rotate=90]
    at (axis cs:-50,2.5) {Real gas region};

  % Annotation
  \node[align=center, font=\small, codexblue]
    at (axis cs:250,0.5)
    {\shortstack{All lines extrapolate to $V=0$\\at $T = -273\textdegree\text{C} = 0\,\text{K}$}};

\end{axis}
\end{tikzpicture}
\end{center}
\small All $V$-$T$ lines, when extrapolated, converge at
$T = 0\ \text{K}$ ($-273°\text{C}$) --- absolute zero,
where the gas would have zero volume. This is how
absolute zero was first predicted from gas behaviour.
\end{diagbox}

\begin{examplebox}[14.2]
\stepcounter{workedexample}
A gas occupies $3.0\ \text{L}$ at $27°\text{C}$
at constant pressure.\\
(a) Find its volume at $127°\text{C}$.\\
(b) Find the temperature at which its volume
would be $1.0\ \text{L}$.
\end{examplebox}

\begin{solutionbox}
Convert temperatures to Kelvin:
$T_1 = 27 + 273 = 300\ \text{K}$

\textbf{(a)} $T_2 = 127 + 273 = 400\ \text{K}$
\[
  V_2 = V_1 \times \frac{T_2}{T_1}
  = 3.0 \times \frac{400}{300} = 4.0\ \text{L}
\]

\textbf{(b)}
\[
  T_2 = T_1 \times \frac{V_2}{V_1}
  = 300 \times \frac{1.0}{3.0} = 100\ \text{K}
  = -173°\text{C}
\]
\end{solutionbox}

% ============================================================
\section{The Pressure Law (Gay-Lussac's Law)}
% ============================================================

\begin{lawbox}[The Pressure Law]
At constant volume, the pressure of a fixed mass of
gas is directly proportional to its absolute
temperature:
\[
  P \propto T
  \quad \text{(constant } V \text{ and mass)}
\]
\[
  \frac{P}{T} = \text{constant}
\]
\[
  \boxed{\frac{P_1}{T_1} = \frac{P_2}{T_2}}
  \quad \text{(at constant } V\text{)}
\]
Again, $T$ must be in Kelvin.
\end{lawbox}

\begin{examplebox}[14.3]
\stepcounter{workedexample}
A car tyre has a pressure of $220\ \text{kPa}$
(gauge) at $15°\text{C}$ in the morning. After a
long drive, the tyre temperature rises to $55°\text{C}$.
Atmospheric pressure is $100\ \text{kPa}$.\\
(a) Find the new pressure in the tyre.\\
(b) Find the new gauge pressure.
\end{examplebox}

\begin{solutionbox}
The Pressure Law uses \textit{absolute} (total)
pressure, not gauge pressure.

Absolute pressure: $P_1 = 220 + 100 = 320\ \text{kPa}$\\
$T_1 = 15 + 273 = 288\ \text{K}$,
$T_2 = 55 + 273 = 328\ \text{K}$

\textbf{(a)} New absolute pressure:
\[
  P_2 = P_1 \times \frac{T_2}{T_1}
  = 320 \times \frac{328}{288}
  = 364\ \text{kPa}
\]

\textbf{(b)} New gauge pressure:
$364 - 100 = 264\ \text{kPa}$

The pressure has risen from $220$ to $264\ \text{kPa}$
gauge --- a $20\%$ increase. Tyre manufacturers
recommend checking tyre pressure when \textit{cold}
for this reason. Letting air out of hot tyres to
reduce pressure is dangerous because the pressure
will drop further when the tyre cools.
\end{solutionbox}

% ============================================================
\section{The General Gas Law and Ideal Gas Equation}
% ============================================================

\begin{lawbox}[The General Gas Law]
Combining Boyle's, Charles' and Pressure Laws:
\[
  \boxed{\frac{P_1 V_1}{T_1} = \frac{P_2 V_2}{T_2}}
  \quad \text{(fixed mass of ideal gas)}
\]
This holds whenever any combination of $P$, $V$
and $T$ changes simultaneously.
\end{lawbox}

\begin{lawbox}[The Ideal Gas Equation]
For $n$ moles of ideal gas:
\[
  \boxed{PV = nRT}
\]
where:
\begin{itemize}[noitemsep]
  \item $P$ = pressure (Pa)
  \item $V$ = volume (m$^3$)
  \item $n$ = amount of gas (moles)
  \item $R = 8.314\ \text{J\,mol}^{-1}\text{K}^{-1}$
  (universal gas constant)
  \item $T$ = absolute temperature (K)
\end{itemize}
For $N$ molecules (using Boltzmann's constant
$k_B = R/N_A = 1.38\times10^{-23}\ \text{J\,K}^{-1}$):
\[
  PV = Nk_BT
\]
At STP ($T = 273\ \text{K}$, $P = 1.013\times10^5\ \text{Pa}$):
1 mole of ideal gas occupies $22.4\ \text{L}$ (molar volume).
\end{lawbox}

\begin{examplebox}[14.4]
\stepcounter{workedexample}
A sealed container holds $0.50\ \text{mol}$ of
nitrogen gas at $300\ \text{K}$ and
$1.5\times10^5\ \text{Pa}$.\\
(a) Calculate the volume of the container.\\
(b) The gas is heated at constant volume to
$450\ \text{K}$. Find the new pressure.\\
(c) Find the number of molecules in the container.\\
($R = 8.314\ \text{J\,mol}^{-1}\text{K}^{-1}$,
$N_A = 6.02\times10^{23}\ \text{mol}^{-1}$)
\end{examplebox}

\begin{solutionbox}
\textbf{(a)} Volume:
\[
  V = \frac{nRT}{P}
  = \frac{0.50 \times 8.314 \times 300}
    {1.5\times10^5}
  = \frac{1247.1}{1.5\times10^5}
  = 8.31\times10^{-3}\ \text{m}^3
  = 8.31\ \text{L}
\]

\textbf{(b)} Constant volume, so $P/T = $ const:
\[
  P_2 = P_1 \times \frac{T_2}{T_1}
  = 1.5\times10^5 \times \frac{450}{300}
  = 2.25\times10^5\ \text{Pa}
\]

\textbf{(c)} Number of molecules:
\[
  N = n \times N_A = 0.50 \times 6.02\times10^{23}
  = 3.01\times10^{23}\ \text{molecules}
\]
\end{solutionbox}

% ============================================================
\section{Dalton's Law of Partial Pressures}
% ============================================================

\begin{lawbox}[Dalton's Law of Partial Pressures]
The total pressure of a mixture of non-reacting
gases equals the sum of the partial pressures of
each component gas:
\[
  P_{total} = P_1 + P_2 + P_3 + \cdots
\]
The \textbf{partial pressure} of a gas is the
pressure it would exert if it alone occupied the
entire volume at the same temperature.
\[
  P_i = \frac{n_i RT}{V}
\]
where $n_i$ is the number of moles of gas $i$.
\end{lawbox}

\begin{examplebox}[14.5]
\stepcounter{workedexample}
A container holds $0.30\ \text{mol}$ of oxygen and
$0.70\ \text{mol}$ of nitrogen at $300\ \text{K}$
in a volume of $10.0\ \text{L}$.\\
(a) Calculate the partial pressure of each gas.\\
(b) Calculate the total pressure.\\
($R = 8.314\ \text{J\,mol}^{-1}\text{K}^{-1}$)
\end{examplebox}

\begin{solutionbox}
\textbf{(a)} Partial pressures:
\[
  P_{O_2} = \frac{n_{O_2}RT}{V}
  = \frac{0.30 \times 8.314 \times 300}
    {0.010}
  = \frac{748.3}{0.010}
  = 74\,830\ \text{Pa} \approx 74.8\ \text{kPa}
\]
\[
  P_{N_2} = \frac{0.70 \times 8.314 \times 300}
    {0.010}
  = 174\,594\ \text{Pa} \approx 174.6\ \text{kPa}
\]

\textbf{(b)} Total pressure:
\[
  P_{total} = 74.8 + 174.6 = 249.4\ \text{kPa}
\]
\end{solutionbox}

% ============================================================
\section{Kinetic Theory of Gases}
% ============================================================

\begin{defbox}[Assumptions of the Kinetic Theory of Ideal Gases]
\begin{enumerate}[noitemsep]
  \item A gas consists of a very large number of
  molecules ($N \sim 10^{23}$)
  \item The molecules are in continuous, rapid,
  random motion
  \item The total volume of the molecules is
  negligible compared to the volume of the container
  \item Collisions between molecules and with the
  walls are perfectly elastic (no kinetic energy lost)
  \item There are no intermolecular forces except
  during collisions (molecules travel in straight
  lines between collisions)
  \item The time of a collision is negligible
  compared to the time between collisions
  \item All molecules obey Newton's laws of motion
\end{enumerate}
\end{defbox}

\begin{diagbox}[Kinetic Model: Gas Molecules in a Container]
\begin{center}
\begin{tikzpicture}[scale=1.0, font=\small]

  % Container
  \fill[gray!10] (0,0) rectangle (7,5);
  \draw[very thick] (0,0) rectangle (7,5);

  % Molecules (random positions)
  \foreach \x/\y in {
    0.8/4.2, 1.5/2.1, 2.3/4.5, 3.1/1.3,
    3.8/3.7, 4.5/0.8, 5.2/2.8, 6.1/1.6,
    1.1/0.6, 2.7/2.9, 4.9/4.1, 6.4/3.5,
    0.4/3.0, 5.8/0.4, 3.4/0.3} {
    \fill[codexblue] (\x,\y) circle (3pt);
  }

  % Velocity arrows on selected molecules
  \draw[->, thick, codexred]
    (1.5,2.1) -- (2.3,2.7);
  \draw[->, thick, codexred]
    (3.1,1.3) -- (2.4,0.7);
  \draw[->, thick, codexred]
    (3.8,3.7) -- (4.6,4.3);
  \draw[->, thick, codexred]
    (5.2,2.8) -- (4.3,2.2);
  \draw[->, thick, codexred]
    (6.1,1.6) -- (6.5,0.9);
  \draw[->, thick, codexred]
    (4.5,0.8) -- (5.3,0.3);

  % Collision with wall
  \fill[codexgold] (6.8,3.5) circle (4pt);
  \draw[->, thick, codexgold]
    (6.0,3.5) -- (6.7,3.5);
  \draw[->, thick, codexgold, dashed]
    (6.8,3.5) -- (6.0,3.5);
  \node[right, codexgold, font=\tiny] at (4.8,3.5)
    {collision};

  % Force exerted on wall
  \foreach \y in {0.5,1.5,2.5,3.5,4.5} {
    \draw[->, thick, codexgreen!70]
      (7.0,\y) -- (7.5,\y);
  }
  \node[right, codexgreen, font=\small] at (7.5,2.5)
    {\shortstack{Net force on wall = Pressure}};

  % Temperature label
  \node[align=center, font=\small] at (3.5,-1.2)
    {Higher $T$ $\implies$ faster molecules $\implies$\\
    harder, more frequent wall collisions $\implies$\\
    higher $P$};

  % Labels
  \node[codexblue, font=\tiny] at (0.5,4.8)
    {Molecules};
  \node[codexred, font=\tiny] at (3.5,2.0)
    {\shortstack{Random\\velocities}};

\end{tikzpicture}
\end{center}
\end{diagbox}

\subsection{Kinetic Theory Explanation of the Gas Laws}

The kinetic theory explains why the gas laws work:

\textbf{Boyle's Law} (constant $T$): If we halve the
volume, each molecule travels half the distance before
hitting a wall --- collisions with the wall are twice
as frequent. Double the frequency means double the
pressure. $P \propto 1/V$.

\textbf{Charles' Law} (constant $P$): Raising the
temperature increases molecular speeds. More frequent
and harder collisions $\implies$ pressure rises. To
keep pressure constant, we must increase the volume
(fewer molecules per unit volume, lower collision
frequency). $V \propto T$.

\textbf{Pressure Law} (constant $V$): Raising temperature
increases molecular speeds and hence the force and
frequency of wall collisions. $P \propto T$.

% ============================================================
\section{Kinetic Theory: Connecting Molecules to
Thermodynamics}
% ============================================================

\subsection{Deriving \texorpdfstring{$PV = \frac{1}{3}Nm\langle c^2\rangle$}{PV = (1/3) N m <c\textasciicircum{}2>} }

This is the Level B content. We derive the pressure
of a gas from first principles --- from Newton's laws
applied to individual molecules.

\begin{processbox}[Pressure from First Principles (Level B)]
Consider $N$ molecules each of mass $m$ in a cubic
box of side $L$ (volume $V = L^3$).

Focus on one molecule moving with velocity component
$u_x$ in the $x$-direction toward a wall.

\textbf{Step 1:} Change in momentum per collision
with the wall:
$\Delta p = 2mu_x$ (elastic collision reverses $x$-velocity)

\textbf{Step 2:} Time between successive collisions
with the same wall:
$\Delta t = 2L/u_x$ (molecule travels to opposite wall and back)

\textbf{Step 3:} Force exerted by one molecule on the wall:
\[
  F_{one} = \frac{\Delta p}{\Delta t}
  = \frac{2mu_x}{2L/u_x} = \frac{mu_x^2}{L}
\]

\textbf{Step 4:} Total force from all $N$ molecules:
\[
  F_{total} = \frac{m}{L}\sum_{i=1}^N u_{xi}^2
  = \frac{Nm}{L}\langle u_x^2\rangle
\]
where $\langle u_x^2\rangle$ is the mean square
$x$-velocity.

\textbf{Step 5:} By symmetry of random motion:
$\langle u_x^2\rangle = \langle u_y^2\rangle
= \langle u_z^2\rangle$

Since $\langle c^2\rangle = \langle u_x^2\rangle
+ \langle u_y^2\rangle + \langle u_z^2\rangle$:
\[
  \langle u_x^2\rangle = \frac{1}{3}\langle c^2\rangle
\]

\textbf{Step 6:} Pressure (force per unit area,
wall area $= L^2$):
\[
  P = \frac{F}{A} = \frac{Nm\langle c^2\rangle/3L}{L^2}
  = \frac{Nm\langle c^2\rangle}{3L^3}
\]
\[
  \boxed{PV = \frac{1}{3}Nm\langle c^2\rangle}
\]
where $\langle c^2\rangle$ is the \textbf{mean square
speed} of the molecules, and $\sqrt{\langle c^2\rangle}$
is the \textbf{root mean square (rms) speed} $c_{rms}$.
\end{processbox}

\subsection{Connecting to Temperature}

Comparing $PV = \frac{1}{3}Nm\langle c^2\rangle$
with the ideal gas law $PV = Nk_BT$:
\[
  \frac{1}{3}Nm\langle c^2\rangle = Nk_BT
\]
\[
  \frac{1}{2}m\langle c^2\rangle = \frac{3}{2}k_BT
\]

\begin{keybox}[Average Molecular Kinetic Energy and Temperature]
\[
  \overline{KE} = \frac{1}{2}m\langle c^2\rangle
  = \frac{3}{2}k_BT
\]
The average translational kinetic energy of a
molecule depends \textbf{only on temperature}.
It is independent of the type of gas, the mass of
the molecule, or the pressure.

This gives the rms speed:
\[
  c_{rms} = \sqrt{\langle c^2\rangle}
  = \sqrt{\frac{3k_BT}{m}}
  = \sqrt{\frac{3RT}{M}}
\]
where $M$ is the molar mass (kg\,mol$^{-1}$).
\end{keybox}

\begin{diagbox}[Maxwell-Boltzmann Speed Distribution]
\begin{center}
\begin{tikzpicture}
\begin{axis}[width=11cm, height=6cm,
  xlabel={Molecular speed $c$ (m\,s$^{-1}$)},
  ylabel={Fraction of molecules},
  xmin=0, xmax=1400, ymin=0, ymax=0.0035,
  grid=major, grid style={gray!15},
  ytick=\empty,
  xtick={0,200,400,600,800,1000,1200},
  legend pos=north east]
  % Low T distribution (peak verified at x=sqrt(80000)=283 m/s)
  \addplot[codexblue, very thick, domain=0:1200, samples=150]
    {1200 * x^2 * exp(-x^2/80000) * 3e-9};
  \addlegendentry{Low $T$}
  % High T distribution (peak verified at x=sqrt(160000)=400 m/s)
  \addplot[codexred, very thick, domain=0:1400, samples=150]
    {600 * x^2 * exp(-x^2/160000) * 3e-9};
  \addlegendentry{High $T$}
  % Low T c_mp: dashed blue line + rotated label on left side
  \draw[dashed, codexblue, thick] (axis cs:283,0) -- (axis cs:283,0.00285);
  \node[codexblue, font=\tiny, rotate=90, anchor=south] at (axis cs:255,0.00145) {$c_{mp}$ (low $T$)};
  % High T c_mp: dashed red (no longer duplicated with a dotted line on top)
  \draw[dashed, codexred, thick] (axis cs:400,0) -- (axis cs:400,0.00185);
  \node[codexred, font=\tiny, rotate=90, anchor=south] at (axis cs:378,0.001) {$c_{mp}$ (high $T$)};
  % c_mean: dotted gold at 450
  \draw[dotted, codexgold, line width=1.2pt] (axis cs:450,0) -- (axis cs:450,0.00175);
  \node[codexgold, font=\tiny, rotate=90, anchor=south] at (axis cs:468,0.0009) {$\bar{c}$ (mean)};
  % c_rms: dash-dot green at 490 (distinct style from dotted and dashed)
  \draw[dash dot, codexgreen, line width=1.2pt] (axis cs:490,0) -- (axis cs:490,0.00155);
  \node[codexgreen, font=\tiny, rotate=90, anchor=south] at (axis cs:508,0.0008) {$c_{rms}$};
  % Area = 1 note (positioned left of the legend box)
  \node[font=\tiny, gray, align=center] at (axis cs:1100,0.0028)
    {Total area\\under each\\curve $= 1$};
  % Tail annotation (anchor and position kept inside the axis frame)
  \draw[->, gray, thick] (axis cs:850,0.0004) -- (axis cs:1000,0.0004);
  \node[gray, font=\tiny, align=left, anchor=west] at (axis cs:1010,0.0004)
    {More high-speed\\molecules at high $T$};
\end{axis}
\end{tikzpicture}
\end{center}
\small \textbf{Key:} The Maxwell-Boltzmann distribution shows the spread of molecular speeds.
At higher $T$: the peak shifts right (faster average), the curve broadens, and importantly,
more molecules have very high speeds (the tail extends further). This is why reactions speed up
with temperature --- more molecules can overcome the activation energy.
\end{diagbox}

\begin{examplebox}[14.6]
\stepcounter{workedexample}
Calculate the rms speed of nitrogen molecules
(molar mass $28\ \text{g\,mol}^{-1}$) at $300\ \text{K}$.\\
($R = 8.314\ \text{J\,mol}^{-1}\text{K}^{-1}$)
\end{examplebox}

\begin{solutionbox}
$M = 28\times10^{-3}\ \text{kg\,mol}^{-1}$

\[
  c_{rms} = \sqrt{\frac{3RT}{M}}
  = \sqrt{\frac{3 \times 8.314 \times 300}
    {28\times10^{-3}}}
  = \sqrt{\frac{7482.6}{0.028}}
  = \sqrt{267\,236}
  = 517\ \text{m\,s}^{-1}
\]

At room temperature, nitrogen molecules move at
about $517\ \text{m\,s}^{-1}$ --- roughly $1.5$ times
the speed of sound. For comparison, oxygen ($M = 32$)
has rms speed $\sqrt{28/32} = 0.935$ times that of
nitrogen, about $483\ \text{m\,s}^{-1}$.
\end{solutionbox}

% ============================================================
\section{Ideal vs Real Gases}
% ============================================================

\begin{defbox}[Ideal vs Real Gas]
An \textbf{ideal gas} perfectly obeys $PV = nRT$.
It has:
\begin{itemize}[noitemsep]
  \item Zero molecular volume
  \item No intermolecular forces (except during
  collisions)
\end{itemize}

A \textbf{real gas} deviates from ideal behaviour:
\begin{itemize}[noitemsep]
  \item Molecules have finite volume (important at
  high pressure, when molecules are close together)
  \item Intermolecular attractions exist (important
  at low temperature, when molecules move slowly)
\end{itemize}

Real gases behave most ideally at:
\begin{itemize}[noitemsep]
  \item \textbf{Low pressure} (molecules far apart)
  \item \textbf{High temperature} (molecules moving
  fast, intermolecular forces less significant)
\end{itemize}
\end{defbox}

\begin{diagbox}[$PV/nRT$ vs $P$ for Ideal and Real Gases]
\begin{center}
\begin{tikzpicture}
\begin{axis}[width=10cm, height=6cm,
  xlabel={Pressure $P$ (atm)},
  ylabel={$PV/nRT$},
  xmin=0, xmax=600, ymin=0.6, ymax=1.4,
  grid=major, grid style={gray!15},
  ytick={0.6,0.7,0.8,0.9,1.0,1.1,1.2,1.3,1.4},
  legend style={at={(1.02,1.0)},anchor=north west,draw=gray!40,fill=white}]
  % Ideal gas (horizontal line at 1)
  \addplot[black, very thick, dashed, domain=0:600]{1.0};
  \addlegendentry{Ideal gas}
  % Hydrogen (rises above 1: volume effect dominates)
  \addplot[codexblue, very thick, domain=0:600, samples=50]
    {1.0 + 0.0003*x};
  \addlegendentry{H$_2$ (real)}
  % Nitrogen (dips below then rises)
  \addplot[codexred, very thick, domain=0:600, samples=80]
    {1.0 - 0.0005*x + 0.0000018*x^2};
  \addlegendentry{N$_2$ (real)}
  % CO2 (larger dip: stronger attractions)
  \addplot[codexgreen, very thick, domain=0:600, samples=80]
    {1.0 - 0.0015*x + 0.0000040*x^2};
  \addlegendentry{CO$_2$ (real)}
  % Ideal line label: anchor=east keeps it right-aligned fully inside the frame
  \node[font=\small, black, anchor=east] at (axis cs:597,1.03)
    {Ideal ($PV/nRT = 1$)};
  % Intermolecular attractions annotation
  \draw[->, gray] (axis cs:200,0.72)
    -- (axis cs:150,0.82)
    node[left, font=\tiny]{\shortstack{Intermolecular\\attractions}};
  % Finite volume: label in open space, arrow points to H2 curve
  \node[gray, font=\tiny, align=center] at (axis cs:470,1.32)
    {Finite\\volume};
  \draw[->, gray] (axis cs:455,1.27) -- (axis cs:425,1.127);
\end{axis}
\end{tikzpicture}
\end{center}
\small For an ideal gas, $PV/nRT = 1$ at all pressures.
Real gases deviate: at low pressure, attractions
dominate (below 1); at high pressure, finite molecular
volume dominates (above 1).
\end{diagbox}

% ============================================================
\section{Chapter Summary}
% ============================================================

\begin{summarybox}
\textbf{Boyle's Law:} $P_1V_1 = P_2V_2$
(constant $T$, $P \propto 1/V$)

\textbf{Charles' Law:} $V_1/T_1 = V_2/T_2$
(constant $P$, $V \propto T$)

\textbf{Pressure Law:} $P_1/T_1 = P_2/T_2$
(constant $V$, $P \propto T$)

\textbf{General gas law:}
$P_1V_1/T_1 = P_2V_2/T_2$

\textbf{Ideal gas equation:}
$PV = nRT = Nk_BT$

$R = 8.314\ \text{J\,mol}^{-1}\text{K}^{-1}$,
$k_B = 1.38\times10^{-23}\ \text{J\,K}^{-1}$

\textbf{Dalton's Law:}
$P_{total} = P_1 + P_2 + \cdots$

\textbf{Kinetic theory result:}
$PV = \frac{1}{3}Nm\langle c^2\rangle$

\textbf{Molecular KE and temperature:}
$\frac{1}{2}m\langle c^2\rangle = \frac{3}{2}k_BT$

\textbf{RMS speed:}
$c_{rms} = \sqrt{3RT/M} = \sqrt{3k_BT/m}$

\textbf{Ideal gas:} $T$ always in Kelvin.
Real gases deviate at high $P$ and low $T$.

\textbf{STP:} 1 mol ideal gas occupies 22.4 L.
\end{summarybox}

% ============================================================
\newpage
\begin{exercisebox}[14]

\subsection*{Section A: Multiple Choice}

\textbf{A1.} A gas is compressed from $8.0\ \text{L}$
at $100\ \text{kPa}$ to $2.0\ \text{L}$ at constant
temperature. The new pressure is:

\mcoption{A}{$25\ \text{kPa}$}
\mcoption{B}{$200\ \text{kPa}$}
\mcoption{C}{$400\ \text{kPa}$}
\mcoption{D}{$800\ \text{kPa}$}
\ansref{Ch14-A1}

\vspace{0.4cm}

\textbf{A2.} A gas at $27°\text{C}$ is heated at
constant pressure until its volume doubles. The
new temperature is:

\mcoption{A}{$54°\text{C}$}
\mcoption{B}{$300°\text{C}$}
\mcoption{C}{$327°\text{C}$}
\mcoption{D}{$600°\text{C}$}
\ansref{Ch14-A2}

\vspace{0.4cm}

\textbf{A3.} Which graph correctly shows Boyle's
Law?

\mcoption{A}{$P$ vs $V$: straight line through origin}
\mcoption{B}{$P$ vs $1/V$: straight line through origin}
\mcoption{C}{$P$ vs $V$: horizontal line}
\mcoption{D}{$P$ vs $T$: straight line through origin}
\ansref{Ch14-A3}

\vspace{0.4cm}

\textbf{A4.} According to kinetic theory, the
average translational kinetic energy of a gas
molecule depends on:

\mcoption{A}{The type of gas and its mass only}
\mcoption{B}{The pressure of the gas only}
\mcoption{C}{The absolute temperature only}
\mcoption{D}{Both the pressure and temperature}
\ansref{Ch14-A4}

\vspace{0.4cm}

\textbf{A5.} The rms speed of molecules in a gas
is doubled. The absolute temperature has:

\mcoption{A}{Doubled}
\mcoption{B}{Halved}
\mcoption{C}{Quadrupled}
\mcoption{D}{Increased by $\sqrt{2}$}
\ansref{Ch14-A5}

\vspace{0.4cm}

\textbf{A6.} $2.0\ \text{mol}$ of an ideal gas
occupy $50\ \text{L}$ at $300\ \text{K}$.
The pressure is ($R = 8.314\ \text{J\,mol}^{-1}\text{K}^{-1}$):

\mcoption{A}{$9.97\ \text{kPa}$}
\mcoption{B}{$99.8\ \text{kPa}$}
\mcoption{C}{$997\ \text{kPa}$}
\mcoption{D}{$49.9\ \text{kPa}$}
\ansref{Ch14-A6}

\vspace{0.4cm}

\textbf{A7.} Dalton's Law states that the total
pressure of a gas mixture equals:

\mcoption{A}{The pressure of the most abundant gas}
\mcoption{B}{The average pressure of all gases}
\mcoption{C}{The sum of the partial pressures}
\mcoption{D}{The product of all partial pressures}
\ansref{Ch14-A7}

\vspace{0.4cm}

\textbf{A8.} An ideal gas differs from a real gas
because an ideal gas has:

\mcoption{A}{Heavier molecules}
\mcoption{B}{Zero molecular volume and no
intermolecular forces}
\mcoption{C}{Molecules moving at the same speed}
\mcoption{D}{All kinetic energy and no potential
energy}
\ansref{Ch14-A8}

\vspace{0.4cm}

\textbf{A9.} At STP, the molar volume of an ideal
gas is approximately:

\mcoption{A}{$2.24\ \text{L}$}
\mcoption{B}{$22.4\ \text{L}$}
\mcoption{C}{$224\ \text{L}$}
\mcoption{D}{$2240\ \text{L}$}
\ansref{Ch14-A9}

\vspace{0.4cm}

\textbf{A10.} The main reason real gases deviate
from ideal gas behaviour at high pressure is that:

\mcoption{A}{Molecules move too fast}
\mcoption{B}{The temperature becomes too high}
\mcoption{C}{The volume of the molecules
themselves becomes significant}
\mcoption{D}{Molecules lose kinetic energy
in collisions}
\ansref{Ch14-A10}

\vspace{0.6cm}

\subsection*{Section B: Theory and Calculations}

\textbf{B1.} A diver uses a $12\ \text{L}$ tank of
compressed air at $200\ \text{atm}$ at the surface
($20°\text{C}$).\\[4pt]
(a) Calculate the volume this air would occupy at
surface pressure ($1.0\ \text{atm}$) at $20°\text{C}$.\\
(b) At a depth of $30\ \text{m}$ in seawater
($\rho = 1025\ \text{kg\,m}^{-3}$), the absolute
pressure is $4.07\ \text{atm}$. The air in the
tank cools to $10°\text{C}$. Calculate the volume
the air would occupy at this depth and temperature.\\
(c) How many times longer can the diver breathe
from the tank at $30\ \text{m}$ compared to at
the surface, if each breath uses the same volume
of air?\\
($g = 10\ \text{m\,s}^{-2}$, $P_{atm} = 1.013\times10^5$ Pa)
\hfill\ansref{Ch14-B1}

\vspace{0.4cm}

\textbf{B2.} An LPG (liquefied petroleum gas)
cylinder has an internal volume of $15\ \text{L}$.
It contains gas at a pressure of $8.0\times10^5\ \text{Pa}$
and temperature $25°\text{C}$. The gas is essentially
propane with molar mass $44\ \text{g\,mol}^{-1}$.
($R = 8.314\ \text{J\,mol}^{-1}\text{K}^{-1}$)\\[4pt]
(a) Calculate the number of moles of propane.\\
(b) Calculate the mass of propane in the cylinder.\\
(c) If the temperature rises to $50°\text{C}$ and
the pressure relief valve opens at
$9.5\times10^5\ \text{Pa}$, does the valve open?\\
(d) The cylinder is left in a fire where the
temperature reaches $300°\text{C}$. Calculate the
pressure if no gas escapes. Comment on the danger.
\hfill\ansref{Ch14-B2}

\vspace{0.4cm}

\textbf{B3.} Use the kinetic theory result
$PV = \frac{1}{3}Nm\langle c^2\rangle$ and the
ideal gas equation to answer the following.\\[4pt]
(a) Show that the average kinetic energy of a
molecule is $\frac{3}{2}k_BT$.\\
(b) Calculate the average KE of a nitrogen
molecule at $300\ \text{K}$ and at $600\ \text{K}$.\\
(c) Calculate the rms speed of nitrogen at both
temperatures ($M_{N_2} = 28\times10^{-3}\ \text{kg\,mol}^{-1}$,
$R = 8.314\ \text{J\,mol}^{-1}\text{K}^{-1}$).\\
(d) The internal energy of an ideal monatomic gas
is $U = \frac{3}{2}nRT$. Calculate the internal
energy of $2.0\ \text{mol}$ of helium at $300\ \text{K}$.
\hfill\ansref{Ch14-B3}

\vspace{0.4cm}

\textbf{B4.} A mixture of gases in a $20\ \text{L}$
container at $300\ \text{K}$ contains:
$1.0\ \text{mol}$ of helium,
$0.5\ \text{mol}$ of neon, and
$1.5\ \text{mol}$ of argon.
($R = 8.314\ \text{J\,mol}^{-1}\text{K}^{-1}$)\\[4pt]
(a) Calculate the partial pressure of each gas.\\
(b) Calculate the total pressure using Dalton's
Law.\\
(c) Verify using $PV = n_{total}RT$.\\
(d) If the argon is removed by a selective
membrane, what is the new total pressure?\\
(e) Helium has molar mass $4\ \text{g\,mol}^{-1}$
and argon has $40\ \text{g\,mol}^{-1}$. At the
same temperature, calculate the ratio of their
rms speeds.
\hfill\ansref{Ch14-B4}

\vspace{0.4cm}

\textbf{B5.} (a) Starting from the ideal gas
equation $PV = nRT$, show that the density $\rho$
of an ideal gas is given by:
\[
  \rho = \frac{PM}{RT}
\]
where $M$ is the molar mass.\\[4pt]
(b) Use this result to calculate the density of
dry air (approximate $M_{air} = 29\ \text{g\,mol}^{-1}$)
at $20°\text{C}$ and 1 atm ($1.013\times10^5\ \text{Pa}$).\\
(c) Hot air balloons work because heated air is
less dense than the surrounding cool air. At what
temperature must air be heated for the density to
be $20\%$ less than at $20°\text{C}$ (with
pressure constant)?\\
(d) On a hot day in Lagos ($35°\text{C}$,
$P = 1.01\times10^5\ \text{Pa}$), calculate the
density of air and compare with the value at $20°\text{C}$.
Explain why aeroplanes require longer runways to
take off on hot days.
\hfill\ansref{Ch14-B5}

\end{exercisebox}